\section{Preliminaries}\label{sec:prelim}

For $n\ge1$, write $[n]=\{1,\ldots,n\}$.  Let $X\subseteq\R^d$ and let
$\cU=(U_1,\ldots,U_n)$ be a finite family of subsets of $X$.  For
$x\in X$, its membership pattern is
\[
  \patt_{\cU}(x)=\{i\in[n]:x\in U_i\}.
\]
The associated code is
\[
  \code(\cU;X)=\{\patt_{\cU}(x):x\in X\}\subseteq2^{[n]}.
\]
Equivalently, a word $\sigma\subseteq[n]$ occurs when its atom
\[
 A_\sigma(\cU;X)
 =X\cap\bigcap_{i\in\sigma}U_i
       \setminus\bigcup_{j\notin\sigma}U_j
\]
is nonempty.  We call the code \emph{open convex} when $X$ and all $U_i$
are open and convex in a common finite-dimensional affine space.

Throughout the planar theorem, $F\subset\R^2$ is a compact convex polygon.
The relative source code is
\[
 \code_F(\cU)=
 \code(U_1\cap\operatorname{int}F,\ldots,
       U_n\cap\operatorname{int}F;\operatorname{int}F).
\]
A \emph{protected witness} is a selected point
$x\in\operatorname{int}F$ together with its source word
$\patt_{\cU}(x)$.  Preserving it means that its membership word is unchanged
in the output realization.

Let $L$ be an affine line.  The \emph{open trace} of $U_i$ on $L$ relative
to $F$ is
\[
  L\cap U_i\cap F,
\]
and the corresponding \emph{closed trace} is
$L\cap\overline{U_i}\cap F$.  The ambient factor $F$ is closed in these
expressions.  Thus an open source interval can become half-open or closed
after clipping by $F$.

A compact convex polygon is permitted to be empty when the corresponding
neuron is trivial.  In the one- and zero-dimensional relative versions,
``polygon'' means respectively a compact interval or a point, with the same
empty-set convention.

\begin{remark}[Relative and ambient interiors]
The output field associated with a two-dimensional polygon $P_i$ is its
Euclidean interior.  On a segment or a point carrier, all interiors are
understood in the relative affine hull.  A closed tangent contact may be
retained even though it does not belong to the strict open trace.
\end{remark}
